% part: intuitionistic-logic
% chapter: semantics
% section: introduction

\documentclass[../../../include/open-logic-section]{subfiles}

\begin{document}

\olfileid[zh]{int}{sem}{int}

\olsection{导论}

没有任何逻辑能在缺失语义的情况下得到令人满意的刻画，而直觉主义逻辑也不例外。对于经典逻辑而言，基于!!{valuation}的语义是典范的，而直觉主义逻辑则有若干相互竞争的语义。它们之中没有一个是完全令人满意的，因为没有一个能在直觉主义上可接受地说明联结词的意义。

基于!!{relational model}的语义与模态逻辑的语义相似，或许是最为流行的一种。在这种语义中，!!{propositional variable}被指派给世界，而这些世界由可及关系相联系。该关系总是偏序，即它是自反的、反对称的且传递的。

直观上，你可以把这些世界想成知识状态或「证据情形」。一个状态~$w'$ 从~$w$ 可及，当且仅当，就我们目前所知，$w'$ 是一个可能的（未来的）知识状态，即与~$w$ 处已知内容相容的状态。一旦一个命题被知道，它就不会变成未知的，即无论何时 $!A$ 在~$w$ 处已知且~$Rww'$，$!A$ 在~$w'$ 处也已知。所以「知识」相对于可及关系是单调的。

如果我们像在认知逻辑中那样把「$!A$ 已知」定义为「在所有认知选项中都为真」，那么 $!A \land !B$ 在~$w$ 处已知，当且仅当在所有认知选项中 $!A$ 与~$!B$ 都已知。但因知识是单调的且 $R$ 是自反的，这意味着 $!A \land !B$ 在 $w$ 处已知，当且仅当 $!A$ 与 $!B$ 在~$w$ 处都已知。出于同样的理由，$!A \lor !B$ 在 $w$ 处已知，当且仅当其中至少一个是已知的。所以对 $\land$ 和 $\lor$ 而言，联结词的真值条件与经典逻辑中的一致。

然而，条件句的真值条件则不同于经典逻辑。$!A \lif !B$ 在~$w$ 处已知，当且仅当不存在满足 $Rww'$ 的~$w'$ 使得 $!A$ 已知而 $!B$ 尚未已知。这不同于「$!A$ 未知或 $!B$~在~$w$ 处已知」这一条件。因为若我们在~$w$ 处既不知 $!A$ 也不知 $!B$，则可能存在一个满足 $Rww'$ 的未来的知识状态~$w'$，使得在 $w'$ 处 $!A$ 已知而 $!B$ 尚未被知道。

我们只有在不存在任何可能的未来知识状态下我们已知~$!A$ 时，才知道 $\lnot !A$。此处的想法是：若 $!A$ 本可以是可知的，则在某个可能的未来知识状态下~$!A$ 成为已知的。由于我们不能知道 $\lfalse$，于是在那个未来的知识状态中，我们会知道 $!A$ 但不知道~$\lfalse$。

在这种解释下，排中律会失效。因为存在一些我们尚未知道、但将来可能会知道的 $!A$。对于这样的!!a{formula}~$!A$，$!A$ 与~$\lnot !A$ 都未知，从而 $!A \lor \lnot !A$ 并不被知道。但我们确实知道，例如 $\lnot(!A \land \lnot !A)$。因为不存在一个我们既知道 $!A$ 又知道 $\lnot !A$ 的未来状态，而且无论我们是否知道~$!A$ 或 $\lnot !A$，这一点都是知道的。

!!^{relational model} 并非直觉主义逻辑唯一可用的语义。拓扑语义是另一种：在这里，命题被解释为拓扑空间中的开集，联结词被解释为这些集合上的运算（例如，$\land$ 对应于交集）。

\end{document}
